(i) By Fact \claimref{standard-model.imported-clebsch-structure-seesaw}, $M_{u} - M_{\nu} = 4w_{u}f + (c_{u} - c_{\nu})h$ and $M_{d} - M_{e} = 4w_{d}f + (c_{d} - c_{e})h$: each $10$ enters the two members of a difference with the same vev and cancels, and the $126$ enters with Clebsches $1$ and $-3$. By Fact \claimref{standard-model.imported-spinor-bilinears}, $16 \otimes 16 = 10_{s} \oplus 120_{a} \oplus 126_{s}$, so $f^{\mathsf{T}} = f$ and $h^{\mathsf{T}} = -h$ (as family matrices of Yukawa couplings to symmetric, respectively antisymmetric, components). Taking symmetric and antisymmetric parts gives the four statements, and with $M_{R} = v_{R}f$ (the same Fact) the displayed identities follow when $w_{u}, w_{d} \neq 0$.

(ii) The hypothesis is the content of Premise (EW) (\claimref{standard-model.premise-ew-equal-weight}), stated there in a family basis in which $\mathrm{Sym}\,M_{u}$ is diagonal with $(M_{u})_{33} = m_{t}$; the computation needs only the two $33$ entries. The diagonal entries of $h$ vanish, so by (i) $(M_{u} - M_{\nu})_{33} = 4w_{u}f_{33}$, that is $4w_{u}f_{33} = m_{t}$, and $(M_{R})_{33} = v_{R}f_{33} = m_{t}v_{R}/(4w_{u})$.

The check builds complex rational instances (two $10$s, symmetric $f$, antisymmetric $h$, generic complex vevs) and verifies (i) exactly, together with the control that a symmetric $120$ coupling would enter $\mathrm{Sym}(M_{d} - M_{e})$; it then imposes (EW) along a heavy direction $x$ ($x^{\mathsf{T}}M_{\nu}x = 0$, $x^{\mathsf{T}}M_{u}x = m_{t}\,x\cdot x$, solved for one $10$ vev and the up-sector normalisation) and verifies $x^{\mathsf{T}}M_{R}x = v_{R}m_{t}\,x\cdot x/(4w_{u})$.

Numbers (not used in the proof; record level). On the chain with the $(15,2,2)$, $v_{R} = 4.2 \times 10^{10}$ GeV (\claimref{standard-model.what-angle-asks-neutrino-sector}). The store's statement that a top-like Dirac mass would need $f_{J}$ near $4000$ corresponds to $m_{t}(M_{I}) \approx 1.44\sqrt{4000} \approx 91$ GeV, which is used here. Then $f_{33} = m_{t}/(4w_{u}) \approx 23\ \mathrm{GeV}/w_{u}$. Perturbativity $f_{33} \le 4\pi$ gives $|w_{u}| \ge m_{t}/(16\pi) \approx 1.8$ GeV and $(M_{R})_{33} \le 5.3 \times 10^{11}$ GeV; $f_{33} \le \sqrt{4\pi}$ gives $|w_{u}| \ge 6.4$ GeV. If $|w_{u}|$ is bounded by the electroweak vev $v \approx 174$ GeV (this needs the normalisation of $w_{u}$ in the Clebsch formula against the canonically normalised doublet, which is cited from memory and is of order one), then $(M_{R})_{33} \ge m_{t}v_{R}/(4v) \approx 5.5 \times 10^{9}$ GeV. These numbers are not compared here with the neutrino paper's valuation of $M$, which rests on Hypothesis (U), a premise that contradicts (EW).
